본 연구에서는 LLM 기반 멀티 에이전트 시스템(MAS)의 문서 생성 성능을 결정짓는 핵심 요소로서 오케스트레이션 구조의 영향을 정량적으로 분석하였다. 고정된 그래프 기반의 Code 구조, 동적 라우팅 기반의 LLM 구조, 그리고 계층적 제어를 수행하는 Hybrid 구조라는 세 가지 제어 메커니즘을 정의하고, 이를 실제 산업 현장의 정형 보고서 및 창의적 글쓰기 작업에 적용하여 자원 효율성과 생성 품질 간의 상관관계를 검증하였다.

\subsection{연구 결과 요약 및 핵심 발견}

실험 결과, 오케스트레이션 구조의 제어 복잡도가 낮을수록 자원 효율성이 향상되지만, 작업의 성격에 따라 요구되는 최적의 구조가 다름을 확인하였다. 특히 정형 데이터 기반의 보고서 생성 작업에서 Code 구조는 결정론적인 워크플로우를 통해 오버헤드를 최소화하였다. 실측 데이터에 따르면, '금형-사출-사출기-단위-가동품질-통계-분석-보고서' 생성 작업 시 Code 구조는 평균 실행 시간 158.61초, 총 토큰 소모량 21,440개를 기록하며 자원 효율성 측면에서 가장 우수한 성능을 보였다. 품질 측면에서도 LLM-as-Judge 평가 결과 82점의 높은 점수를 기록하였으며, 이는 정형화된 파이프라인만으로도 충분한 품질의 전문 보고서 작성이 가능함을 시사한다.

반면, 제어 흐름이 동적으로 결정되는 LLM 및 Hybrid 구조는 에이전트 간의 상호작용과 상태 분석을 위한 추가적인 추론 단계가 필요하므로, Code 구조 대비 실행 시간과 토큰 소모량이 증가하는 트레이드-오프(Trade-off) 관계가 나타났다. 그러나 이는 복잡한 문맥적 판단이 필요하거나 창의적인 구성이 요구되는 작업에서 생성물의 유연성과 완성도를 높이는 핵심 기제로 작용함을 확인하였다.

\subsection{최종 결론}

결론적으로, 본 연구는 서론에서 제기한 "오케스트레이션 구조가 멀티 에이전트 시스템의 성능과 효율성에 어떠한 영향을 미치는가"라는 문제에 대해, 제어 메커니즘의 결정론적 성격과 동적 성격의 정도에 따라 자원 소모량과 품질의 균형점이 달라진다는 명확한 답을 제시하였다.

본 연구의 결과를 바탕으로 작업 성격에 따른 최적의 오케스트레이션 선택 기준을 다음과 같이 제안한다. 첫째, 산업 보고서와 같이 입력 데이터가 정형적이고 출력 형식이 엄격히 정의된 작업의 경우, 낮은 오버헤드와 예측 가능한 성능을 보장하는 Code 구조가 가장 효율적이다. 둘째, 작업의 목표가 모호하거나 생성 과정에서 지속적인 수정과 피드백 루프가 필수적인 창의적 작업의 경우, 자원 소모가 증가하더라도 LLM 또는 Hybrid 구조를 채택하여 동적 라우팅을 통한 품질 최적화를 도모해야 한다.

본 연구는 멀티 에이전트 시스템 설계 시 단순히 모델의 성능에 의존하는 것이 아니라, 해결하고자 하는 과업의 특성에 맞추어 오케스트레이션 구조를 전략적으로 선택해야 함을 정량적으로 입증하였다는 점에서 학술적 및 실무적 의의를 갖는다. 향후 연구에서는 제안된 세 가지 구조 외에, 강화학습을 통해 최적의 경로를 스스로 학습하는 적응형 오케스트레이션 구조로 확장하여 효율성과 효과성을 동시에 극대화하는 방안을 모색할 필요가 있다.